\documentclass[11pt]{article}

\usepackage[spanish]{babel}
\usepackage[utf8]{inputenc}
\usepackage[T1]{fontenc}
\usepackage{amsmath, amssymb, amsthm, mathtools}
\usepackage{geometry}
\usepackage{xcolor}
\usepackage{tcolorbox}
\usepackage{makecell}

\geometry{letterpaper, margin=2.2cm}

% ==============================================================================
% DEFINICIONES DE ENTORNOS PERSONALIZADOS PARA COMPILACIÓN
% ==============================================================================

\newenvironment{motivacion}{%
  \section*{1. Motivación}%
}{}

\newenvironment{preguntaguia}{%
  \begin{tcolorbox}[colback=cyan!5,colframe=cyan!75!black,title=Pregunta Guía]%
}{%
  \end{tcolorbox}%
}

\newenvironment{teoria}{%
  \section*{2. Definición y Teoría}%
}{}

\newenvironment{definicion}[1]{%
  \begin{tcolorbox}[colback=blue!5,colframe=blue!75!black,title=Definición: #1]%
}{%
  \end{tcolorbox}%
}

\newenvironment{teorema}[1]{%
  \begin{tcolorbox}[colback=green!5,colframe=green!75!black,title=Teorema: #1]%
}{%
  \end{tcolorbox}%
}

\newenvironment{alerta}[1]{%
  \begin{tcolorbox}[colback=red!5,colframe=red!80!black,title=Advertencia Metodológica: #1]%
}{%
  \end{tcolorbox}%
}

\newenvironment{repaso}[1]{%
  \begin{tcolorbox}[colback=orange!5,colframe=orange!80!black,title=Recordatorio de Álgebra Lineal: #1]%
}{%
  \end{tcolorbox}%
}

\newenvironment{procesamiento}[1]{%
  \begin{tcolorbox}[colback=gray!5,colframe=gray!60!black,title=Procedimiento: #1]%
}{%
  \end{tcolorbox}%
}

\newenvironment{aplicacion}{%
  \section*{3. Aplicación y Práctica}%
}{}

\newenvironment{ejemplo}[1]{%
  \begin{tcolorbox}[colback=white,colframe=gray!40,title=Ejemplo: #1]%
}{%
  \end{tcolorbox}%
}

% Comandos para preguntas interactivas
\newenvironment{preguntaalternativas}[1]{\subsection*{Pregunta: #1}\par\vspace{0.1cm}}{\vspace{0.3cm}}
\newcommand{\opcion}[3]{\par\noindent $\bullet$ \textbf{#1} \textit{(#2)} - #3\par\vspace{0.15cm}}

\newenvironment{preguntaverdaderofalso}[1]{\subsection*{Pregunta V/F: #1}\par\vspace{0.1cm}}{\vspace{0.3cm}}
\newcommand{\verdaderofalso}[3]{\par\noindent\textbf{Respuesta correcta: #1} \\ \textit{Si marca Falso:} #2 \\ \textit{Si marca Verdadero:} #3\par\vspace{0.15cm}}

% Entorno Términos Pareados (2 Columnas)
\newenvironment{pareados}[1]{\subsection*{Términos Pareados: #1}}{\vspace{0.3cm}}
\newcommand{\columnaI}[1]{\textbf{Columna 1:}\begin{enumerate}#1\end{enumerate}}
\newcommand{\columnaII}[1]{\textbf{Columna 2:}\begin{itemize}#1\end{itemize}}
\newcommand{\pareo}[2]{\textbf{Cruce #1:} #2\par}

% Entornos para ejercicios
\newenvironment{ejercicios}{%
  \section*{4. Ejercicios Resueltos y Propuestos}%
}{}

\newenvironment{ejercicioresuelto}[2]{\subsection*{Ejercicio Resuelto: #1 \small{[#2]}}}{}

\newcommand{\enunciado}[1]{\textbf{Enunciado:} #1\par\vspace{0.2cm}}
\newcommand{\solucion}[1]{\textbf{Solución:} #1\par\vspace{0.4cm}}

% Entorno Fórmulas
\newenvironment{formulas}{\section*{5. Fórmulas Clave}\begin{description}}{\end{description}}
\newcommand{\formula}[3]{\item[#1:] $#2$ \\ \textit{#3} \vspace{0.2cm}}

% ==============================================================================
% METADATOS TÉCNICOS DE DESPLIEGUE
% ==============================================================================
% ID_CURSO: calculo-multivariable
% NUMERO_UNIDAD: 3
% TITULO_UNIDAD: Diferenciabilidad
% NUMERO_CAPITULO: 3.2
% TITULO_CAPITULO: Diferenciabilidad y Plano Tangente
% ESTADO_BLOQUEADO: false
% ESTADO_COMPLETADO: false
% ==============================================================================

\begin{document}

% ==============================================================================
% PESTAÑA 1: MOTIVACIÓN (contentMotivation)
% ==============================================================================
\begin{motivacion}
  \textbf{Aproximación Lineal y Diferenciabilidad}

  En el análisis de campos escalares, la existencia de derivadas parciales en un punto no constituye una condición suficiente para garantizar la continuidad de la función.

  Considérese, a modo de ilustración, la función analizada en la sección previa:
  $$ f(x,y) = \begin{cases} \dfrac{xy^2}{x^2 + y^4} & \text{si } (x,y) \neq (0,0) \\[1em] 0 & \text{si } (x,y) = (0,0) \end{cases} $$
  En el origen $(0,0)$, ambas derivadas parciales existen y toman el valor cero ($\dfrac{\partial f}{\partial x}(0,0) = 0$ y $\dfrac{\partial f}{\partial y}(0,0) = 0$). Sin embargo, la función no es continua en $(0,0)$, pues al aproximarse al origen a través de trayectorias parabólicas de la forma $x = my^2$, el valor del límite depende del parámetro $m$.

  Esta limitación obedece a que las derivadas parciales evalúan el comportamiento de la función exclusivamente a lo largo de las direcciones ortogonales de los ejes coordenados. Para extender la noción de derivabilidad de una variable real a dimensiones superiores, se requiere exigir una condición analítica más fuerte: la \textbf{diferenciabilidad}. Esta propiedad impone que los incrementos de la función en un entorno del punto puedan ser aproximados mediante un operador lineal, garantizando que el residuo de dicha estimación decaiga a cero más rápidamente que la distancia al punto de evaluación.

  \begin{preguntaguia}
    Dada la existencia de una transformación lineal que aproxima a un campo escalar $f$ en un entorno de un punto, ¿qué relación existe entre los coeficientes de dicha transformación lineal y las derivadas parciales de $f$?
  \end{preguntaguia}
\end{motivacion}


% ==============================================================================
% PESTAÑA 2: DEFINICIÓN Y TEORÍA (contentTheory)
% ==============================================================================
\begin{teoria}
  
  El estudio de la diferenciabilidad en campos escalares se fundamenta en la capacidad de aproximar localmente el comportamiento de una función mediante una transformación lineal.

  % --- DEFINICIÓN: DIFERENCIABILIDAD ---
  \begin{definicion}{Diferenciabilidad y Aproximación Lineal}
    Sea $f \colon D \subseteq \mathbb{R}^2 \to \mathbb{R}$ un campo escalar y sea $(x_0, y_0) \in D$ tal que existe $\delta > 0$ con $B((x_0, y_0), \delta) \subseteq D$. Se dice que $f$ es \textbf{diferenciable} en $(x_0, y_0)$ si existe una \textbf{transformación lineal} $T \colon \mathbb{R}^2 \to \mathbb{R}$, definida por $T(h,k) = Ah + Bk$ con $A, B \in \mathbb{R}$, tal que:
    $$ \lim\limits_{(x,y) \to (x_0, y_0)} \dfrac{f(x,y) - \Big( f(x_0, y_0) + A(x-x_0) + B(y-y_0) \Big)}{\| (x,y) - (x_0, y_0) \|} = 0 $$
  \end{definicion}

  % --- TEOREMA: UNICIDAD Y DEMOSTRACIÓN ---
  \begin{teorema}{Unicidad de la Transformación Lineal y Existencia de Derivadas Parciales}
    Si $f \colon D \subseteq \mathbb{R}^2 \to \mathbb{R}$ es diferenciable en $(x_0, y_0)$, entonces las derivadas parciales de $f$ en dicho punto existen, la transformación lineal $T(h,k) = Ah + Bk$ es \textbf{única} y sus coeficientes están determinados unívocamente por:
    $$ A = \dfrac{\partial f}{\partial x}(x_0, y_0) \quad \text{y} \quad B = \dfrac{\partial f}{\partial y}(x_0, y_0) $$
  \end{teorema}

  \begin{proof}
    Por hipótesis, $f$ es diferenciable en $(x_0, y_0)$, lo que garantiza la existencia de constantes $A, B \in \mathbb{R}$ tales que:
    $$ \lim\limits_{(h,k) \to (0,0)} \dfrac{f(x_0+h, y_0+k) - f(x_0, y_0) - (Ah + Bk)}{\|(h,k)\|} = 0 $$
    Dado que este límite existe y es igual a cero a lo largo de cualquier trayectoria que converja a $(0,0)$, evaluamos el límite a lo largo de las direcciones de los ejes coordenados:

    \begin{enumerate}
      \item \textbf{Aproximación a lo largo del eje $X$ ($k = 0, h \to 0$):}
      Sustituyendo $k = 0$ se obtiene $\|(h,0)\| = |h|$, de modo que:
      $$ \lim\limits_{h \to 0} \dfrac{f(x_0+h, y_0) - f(x_0, y_0) - Ah}{|h|} = 0 $$
      Esto implica analíticamente que:
      $$ \lim\limits_{h \to 0} \left| \dfrac{f(x_0+h, y_0) - f(x_0, y_0)}{h} - A \right| = 0 \implies \lim\limits_{h \to 0} \dfrac{f(x_0+h, y_0) - f(x_0, y_0)}{h} = A $$
      Por definición de derivada parcial, este límite corresponde a $\dfrac{\partial f}{\partial x}(x_0, y_0)$. En consecuencia, $\dfrac{\partial f}{\partial x}(x_0, y_0)$ existe y $A = \dfrac{\partial f}{\partial x}(x_0, y_0)$.

      \item \textbf{Aproximación a lo largo del eje $Y$ ($h = 0, k \to 0$):}
      De manera análoga, al aproximarnos por el eje Y:
      $$ \lim\limits_{k \to 0} \dfrac{f(x_0, y_0+k) - f(x_0, y_0)}{k} = B $$
      Por lo tanto, $\dfrac{\partial f}{\partial y}(x_0, y_0)$ existe y $B = \dfrac{\partial f}{\partial y}(x_0, y_0)$.
    \end{enumerate}
    La transformación lineal es, por tanto, única y sus coeficientes son las derivadas parciales.
  \end{proof}

  % --- DEFINICIÓN: VECTOR GRADIENTE ---
  \begin{definicion}{Vector Gradiente}
    Dada una función $f \colon D \subseteq \mathbb{R}^2 \to \mathbb{R}$ cuyas derivadas parciales de primer orden existen en un punto $(x_0, y_0) \in D$, se define el \textbf{Vector Gradiente} de $f$ en $(x_0, y_0)$, denotado por $\nabla f(x_0, y_0)$, como el vector:
    $$ \nabla f(x_0, y_0) = \left( \dfrac{\partial f}{\partial x}(x_0, y_0), \; \dfrac{\partial f}{\partial y}(x_0, y_0) \right) $$
  \end{definicion}

  % --- DEFINICIÓN: PLANO TANGENTE ---
  \begin{definicion}{Plano Tangente}
    Sea $f \colon D \subseteq \mathbb{R}^2 \to \mathbb{R}$ un campo escalar diferenciable en $(x_0, y_0)$. Se define el \textbf{plano tangente} a la superficie $z = f(x,y)$ en el punto $(x_0, y_0, f(x_0, y_0))$ mediante la ecuación cartesiana:
    $$ z = f(x_0, y_0) + \nabla f(x_0, y_0) \cdot (x - x_0, y - y_0) $$
    O equivalentemente:
    $$ z = f(x_0, y_0) + \left[ \dfrac{\partial f}{\partial x}(x_0, y_0) \right] (x - x_0) + \left[ \dfrac{\partial f}{\partial y}(x_0, y_0) \right] (y - y_0) $$
  \end{definicion}

  % --- PROCEDIMIENTO DE EVALUACIÓN ---
  \begin{procesamiento}{Procedimiento Analítico para Verificar Diferenciabilidad}
    Para determinar rigurosamente la diferenciabilidad de un campo escalar $f$ en un punto $(x_0, y_0)$:
    \begin{enumerate}
      \item Calcule las derivadas parciales en $(x_0, y_0)$. Si la función presenta discontinuidades o está definida a tramos, utilice la definición por límite del cociente incremental.
      \item Formule la aproximación lineal candidata $L(x,y) = f(x_0, y_0) + \dfrac{\partial f}{\partial x}(x_0, y_0)(x-x_0) + \dfrac{\partial f}{\partial y}(x_0, y_0)(y-y_0)$.
      \item Evalúe el límite del residuo de la aproximación:
      $$ \lim\limits_{(x,y) \to (x_0, y_0)} \dfrac{f(x,y) - L(x,y)}{\| (x,y) - (x_0, y_0) \|} $$
      La función es diferenciable en $(x_0, y_0)$ si y solo si dicho límite existe y es igual a cero.
    \end{enumerate}
  \end{procesamiento}

\end{teoria}


% ==============================================================================
% PESTAÑA 3: APLICACIÓN Y PRÁCTICA (contentApplication)
% ==============================================================================
\begin{aplicacion}

  % --- REPASO DE ÁLGEBRA LINEAL ---
  \begin{repaso}{Planos en $\mathbb{R}^3$}
    Para resolver analíticamente problemas geométricos que involucran planos tangentes, es imperativo dominar los siguientes postulados del Álgebra Lineal:
    \begin{enumerate}
      \item \textbf{Vector Normal:} Todo plano se define por un punto $P_0(x_0, y_0, z_0)$ y un vector normal ortogonal a la superficie $\vec{n} = (A, B, C)$. Su ecuación es $A(x-x_0) + B(y-y_0) + C(z-z_0) = 0$.
      \item \textbf{Normal del Plano Tangente:} Al reordenar la ecuación del plano tangente $z - z_0 = f_x(x-x_0) + f_y(y-y_0)$, se deduce inmediatamente que su vector normal es $\vec{n} = \left( \dfrac{\partial f}{\partial x}, \dfrac{\partial f}{\partial y}, -1 \right)$.
      \item \textbf{Condición de Paralelismo:} Dos planos son paralelos si y solo si sus vectores normales son colineales (proporcionales): $\vec{n}_1 = k\vec{n}_2$ para algún $k \in \mathbb{R} \setminus \{0\}$.
      \item \textbf{Tangencia entre Superficies:} Dos superficies $z = f(x,y)$ y $z = g(x,y)$ son tangentes en un punto $P_0$ si coinciden en dicho punto ($f(P_0) = g(P_0)$) y sus vectores gradientes son idénticos ($\nabla f(P_0) = \nabla g(P_0)$), lo que implica que comparten el mismo plano tangente.
      \item \textbf{Ecuación Segmentaria o Simétrica del Plano:} Un plano que interseca los ejes coordenados en $(a,0,0)$, $(0,b,0)$ y $(0,0,c)$ (con $a,b,c \neq 0$) puede escribirse en la forma canónica: $\dfrac{x}{a} + \dfrac{y}{b} + \dfrac{z}{c} = 1$.
    \end{enumerate}
  \end{repaso}
  
  Se evalúa la comprensión profunda de las implicaciones lógicas de la diferenciabilidad y el cálculo algebraico de aproximaciones lineales.

  % --- BATERÍA DE VERDADERO Y FALSO: IMPLICACIONES LÓGICAS ---
  \begin{preguntaverdaderofalso}{Implicación Analítica I: Derivadas Parciales}
    Si un campo escalar $f(x,y)$ posee derivadas parciales bien definidas y finitas en un punto $(x_0,y_0)$, entonces $f$ es diferenciable en dicho punto.
    
    \verdaderofalso{F}
      {¡Correcto! La existencia de derivadas parciales es una condición \textbf{necesaria}, pero no \textbf{suficiente}. Solo aseguran la existencia de tasas de cambio en las direcciones X e Y, pero no garantizan que la superficie completa admita un plano tangente (linealidad local).}
      {Incorrecto. Cuidado con el error clásico. Las derivadas parciales pueden existir y, sin embargo, la función podría incluso no ser continua. Recuerde el análisis patológico del capítulo anterior.}
  \end{preguntaverdaderofalso}

  \begin{preguntaverdaderofalso}{Implicación Analítica II: Existencia del Gradiente}
    Si el vector gradiente $\nabla f(x_0, y_0)$ existe para un campo escalar $f$, entonces queda garantizado matemáticamente que $f$ es diferenciable en $(x_0, y_0)$.
    
    \verdaderofalso{F}
      {¡Excelente! Esta proposición es equivalente a la anterior pero camuflada en notación vectorial. El gradiente es simplemente un vector que agrupa las derivadas parciales. Su existencia no certifica que el límite del residuo de la aproximación decaiga a cero.}
      {Incorrecto. El gradiente existe si existen $\partial f/\partial x$ y $\partial f/\partial y$. Poseer esos coeficientes permite ``escribir'' un plano candidato, pero si $f$ no es diferenciable, dicho plano es un fantasma que no aproxima a la superficie real.}
  \end{preguntaverdaderofalso}

  \begin{preguntaverdaderofalso}{Implicación Analítica III: Condición Necesaria}
    Si un campo escalar $f(x,y)$ es diferenciable en un punto $(x_0,y_0)$, entonces es estrictamente obligatorio que sus derivadas parciales existan en $(x_0,y_0)$.
    
    \verdaderofalso{V}
      {Incorrecto. Repase el Teorema de Unicidad. La diferenciabilidad asegura que la función es ``suave'' localmente, lo cual implica necesariamente que posee derivadas bien definidas en todas las direcciones, incluyendo las parciales.}
      {¡Perfecto! Este es el \textbf{Teorema de Unicidad de la Transformación Lineal}. Si la función admite una aproximación plana (es diferenciable), la matemática exige que las pendientes de ese plano correspondan obligatoriamente a las derivadas parciales.}
  \end{preguntaverdaderofalso}

  % --- PREGUNTA 4: CÁLCULO DE PLANO TANGENTE ---
  \begin{preguntaalternativas}{Cálculo Geométrico del Plano Tangente}
    Determine la ecuación del plano tangente al campo escalar $f(x,y) = x^2 + y^4 + e^{xy}$ en el punto $(x_0, y_0) = (1, 0)$.
    
    \opcion{$z = 2x + y + 2$}{Incorrecto}{Se evaluaron las derivadas parciales correctamente, pero la estructura de la aproximación lineal exige multiplicar por los desplazamientos $(x - x_0)$ y $(y - y_0)$.}
    \opcion{$z = 2x + y$}{Correcto}{Se evalúa $z_0 = f(1,0) = 1 + 0 + 1 = 2$. El gradiente es $\nabla f(x,y) = (2x + y e^{xy}, 4y^3 + x e^{xy})$, de donde $\nabla f(1,0) = (2, 1)$. La ecuación del plano es $z = 2 + 2(x-1) + 1(y-0)$, lo cual simplifica a $z = 2x + y$.}
    \opcion{$z = 2(x-1) + (y-0)$}{Incorrecto}{Esta expresión representa únicamente el incremento diferencial, omitiendo la ordenada de anclaje (altura) $z_0 = f(x_0, y_0) = 2$.}
  \end{preguntaalternativas}

  % --- PREGUNTA 5: TÉRMINOS PAREADOS (GEOMETRÍA DEL PLANO) ---
  \begin{pareados}{Condiciones Geométricas del Plano Tangente}
    Relacione las condiciones analíticas del gradiente con sus correspondientes interpretaciones geométricas para el plano tangente $z = L(x,y)$.

    \columnaI{
      \item $\nabla f(x_0, y_0) = (0,0)$.
      \item $\dfrac{\partial f}{\partial y}(x_0, y_0) = 0$ y $\dfrac{\partial f}{\partial x}(x_0, y_0) \neq 0$.
      \item $\nabla f(x_0, y_0) = (A, B)$ evaluado en un punto donde $f(x_0,y_0) = 0$.
    }

    \columnaII{
      \item El plano tangente intercepta el origen del sistema tridimensional o corta al plano $XY$.
      \item El plano tangente es paralelo al eje $Y$ (no tiene inclinación lateral).
      \item El plano tangente es completamente horizontal (paralelo al plano $XY$).
    }

    \pareo{1-C}{Correcto. Si ambas derivadas parciales son nulas, la ecuación se reduce a $z = f(x_0,y_0)$, lo cual define un plano de altura constante.}
    \pareo{2-B}{Correcto. Si la derivada parcial respecto a $y$ es nula, el plano no experimenta variación al desplazarse en dirección $Y$.}
    \pareo{3-A}{Correcto. Si la altura de evaluación original $z_0$ es cero, la ecuación del plano carece de término independiente y pasa por la coordenada vertical $z=0$.}
  \end{pareados}

\end{aplicacion}


% ==============================================================================
% PESTAÑA 4: EJERCICIOS RESUELTOS Y PROPUESTOS (contentExercises)
% ==============================================================================
\begin{ejercicios}
  A continuación, se desarrollan ejercicios resueltos para asentar la metodología analítica, seguidos de ejercicios propuestos para fomentar la práctica autónoma.

  \subsection*{I. Ejercicios Resueltos}

  % --- EJERCICIO RESUELTO 1 ---
  \begin{ejercicioresuelto}{Estudio Analítico de Diferenciabilidad}{nivel-2}
    \enunciado{
      Considere el campo escalar:
      $$ f(x,y) = \begin{cases} \dfrac{xy}{\sqrt{x^2 + y^2}} & \text{si } (x,y) \neq (0,0) \\[1em] 0 & \text{si } (x,y) = (0,0) \end{cases} $$
      Determine si $f$ es diferenciable en el origen $(0,0)$.
    }
    \solucion{
      \textbf{Paso 1. Cálculo de derivadas parciales en el origen:}
      Al tratarse de una función definida por tramos alrededor de $(0,0)$, evaluamos mediante límites incrementales:
      $$ \dfrac{\partial f}{\partial x}(0,0) = \lim\limits_{h \to 0} \dfrac{f(h,0) - f(0,0)}{h} = \lim\limits_{h \to 0} \dfrac{0}{h} = 0 $$
      Por simetría estructural respecto a $x$ e $y$, se tiene que $\dfrac{\partial f}{\partial y}(0,0) = 0$.
      En consecuencia, $\nabla f(0,0) = (0,0)$ y la aproximación lineal candidata es $L(x,y) = 0$.

      \vspace{0.2cm}
      \textbf{Paso 2. Evaluación del límite de diferenciabilidad:}
      Sustituyendo en la definición formal:
      $$ \lim\limits_{(x,y) \to (0,0)} \dfrac{f(x,y) - L(x,y)}{\| (x,y) \|} = \lim\limits_{(x,y) \to (0,0)} \dfrac{\dfrac{xy}{\sqrt{x^2 + y^2}} - 0}{\sqrt{x^2 + y^2}} = \lim\limits_{(x,y) \to (0,0)} \dfrac{xy}{x^2 + y^2} $$
      Transformando a coordenadas polares ($x=r\cos\theta, y=r\sin\theta$):
      $$ \lim\limits_{r \to 0^+} \dfrac{r^2\cos\theta \sin\theta}{r^2} = \cos\theta \sin\theta $$
      Puesto que el resultado del límite depende de la dirección $\theta$, el límite global no existe. Por consiguiente, \textbf{$f$ no es diferenciable en $(0,0)$}.
    }
  \end{ejercicioresuelto}

  % --- EJERCICIO RESUELTO 2 ---
  \begin{ejercicioresuelto}{Campo Escalar Diferenciable sin Derivadas Parciales Continuas}{nivel-3}
    \enunciado{
      Considere el campo escalar patológico:
      $$ f(x,y) = \begin{cases} \left(x^2 + y^2\right) \sin\left(\dfrac{1}{\sqrt{x^2 + y^2}}\right) & \text{si } (x,y) \neq (0,0) \\[1em] 0 & \text{si } (x,y) = (0,0) \end{cases} $$
      \begin{enumerate}
        \item Calcule las derivadas parciales $\dfrac{\partial f}{\partial x}$ y $\dfrac{\partial f}{\partial y}$ en todo $\mathbb{R}^2$.
        \item Demuestre por definición formal que $f$ es diferenciable en el origen.
        \item Analice la continuidad de las derivadas parciales en el origen.
      \end{enumerate}
    }
    \solucion{
      \textbf{Parte 1: Derivadas parciales en $\mathbb{R}^2$} \\
      Para $(x,y) \neq (0,0)$, aplicando reglas de derivación analítica:
      $$ \dfrac{\partial f}{\partial x} = 2x\sin\left(\dfrac{1}{\sqrt{x^2+y^2}}\right) - \dfrac{x}{\sqrt{x^2+y^2}}\cos\left(\dfrac{1}{\sqrt{x^2+y^2}}\right) $$
      Para $(x,y) = (0,0)$, evaluando el cociente incremental:
      $$ \dfrac{\partial f}{\partial x}(0,0) = \lim\limits_{h \to 0} \dfrac{h^2 \sin\left(\dfrac{1}{|h|}\right) - 0}{h} = \lim\limits_{h \to 0} h \sin\left(\dfrac{1}{|h|}\right) = 0 $$
      Por simetría, $\dfrac{\partial f}{\partial y}(0,0) = 0$.

      \vspace{0.3cm}
      \textbf{Parte 2: Demostración de diferenciabilidad} \\
      Con $\nabla f(0,0) = (0,0)$ y $f(0,0)=0$, se evalúa el límite del residuo:
      $$ \lim\limits_{(x,y) \to (0,0)} \dfrac{(x^2+y^2)\sin\left(\dfrac{1}{\sqrt{x^2+y^2}}\right) - 0}{\sqrt{x^2+y^2}} = \lim\limits_{r \to 0^+} \dfrac{r^2 \sin(1/r)}{r} = \lim\limits_{r \to 0^+} r \sin(1/r) = 0 $$
      Dado que el residuo tiende a cero, \textbf{$f$ es diferenciable en el origen}.

      \vspace{0.3cm}
      \textbf{Parte 3: Continuidad de las derivadas parciales y análisis teórico} \\
      Evaluamos la existencia de $\lim\limits_{(x,y)\to(0,0)} \dfrac{\partial f}{\partial x}$ usando coordenadas polares:
      $$ \lim\limits_{r \to 0^+} \left[ 2r\cos\theta \sin(1/r) - \cos(1/r)\cos\theta \right] $$
      El término $2r\cos\theta \sin(1/r)$ converge a $0$ por acotamiento. Sin embargo, el término $\cos(1/r)\cos\theta$ oscila indefinidamente. En consecuencia, el límite no existe y \textbf{$\dfrac{\partial f}{\partial x}$ no es continua en $(0,0)$}.
    }
  \end{ejercicioresuelto}

  % --- EJERCICIO RESUELTO 3 ---
  

  \vspace{0.5cm}
  \subsection*{II. Ejercicios Propuestos}
  \vspace{0.2cm}

  % --- EJERCICIO PROPUESTO 1 ---
  \textbf{Ejercicio Propuesto 1: Tangencia de Superficies}
  \enunciado{
    Probar rigurosamente que las gráficas de $f(x,y) = x^2 + y^2$ y $g(x,y) = -x^2 - y^2 + xy^3$ son tangentes en el origen $(0,0)$.
  }
  \begin{tcolorbox}[colback=yellow!5,colframe=yellow!60!black,title=Indicación Estratégica]
    Dos superficies son tangentes en un punto en común si comparten exactamente el mismo plano tangente. Verifique primero que $f(0,0)=g(0,0)$ y, a continuación, calcule y compare el vector gradiente de ambas funciones evaluado en el origen. El plano resultante en ambos casos debe ser $z=0$.
  \end{tcolorbox}
  \vspace{0.3cm}

  % --- EJERCICIO PROPUESTO 2 ---
  \textbf{Ejercicio Propuesto 2: Paralelismo y Vectores Normales}
  \enunciado{
    ¿En qué punto $(x_0, y_0, z_0)$ el plano tangente al grafo de la función $f(x,y) = 9 - 4x^2 - y^2$ resulta ser paralelo al plano de ecuación $z = 4y$?
  }
  \begin{tcolorbox}[colback=yellow!5,colframe=yellow!60!black,title=Respuesta y Pista Analítica]
    \textbf{Respuesta:} El único punto de tangencia que cumple la condición es $(0, -2, 5)$. \\
    \textit{Pista:} Recuerde que dos planos son paralelos si sus vectores normales son paralelos (linealmente dependientes). El vector normal genérico del plano tangente a una superficie $z=f(x,y)$ es $\vec{n}_1 = (f_x, f_y, -1)$, y el del plano dado es $\vec{n}_2 = (0, 4, -1)$.
  \end{tcolorbox}
  \vspace{0.3cm}

  % --- EJERCICIO PROPUESTO 3 ---
  \textbf{Ejercicio Propuesto 3: La Constante de las Intersecciones}
  \enunciado{
    Considere la función $f(x,y) = \left( 1-\sqrt{x} -\sqrt{y} \right)^2$. 
    Sean $(a,0,0)$, $(0,b,0)$ y $(0,0,c)$ los puntos de intersección del plano tangente a dicha función con los tres ejes coordenados cartesianos. Demuestre algebraicamente que la suma $a+b+c$ es independiente del punto de tangencia escogido.
  }
  \begin{tcolorbox}[colback=yellow!5,colframe=yellow!60!black,title=Indicación Estratégica]
    Construya la ecuación genérica del plano tangente en un punto arbitrario $(x_0, y_0, z_0)$. Luego, encuentre el intercepto con el eje $X$ evaluando $y=0, z=0$ (obteniendo el valor de $a$), y repita el proceso análogo para $b$ y $c$. Si el álgebra es correcta, al sumar las expresiones resultantes, las variables $x_0$ e $y_0$ se cancelarán, demostrando que $a+b+c = 1$.
  \end{tcolorbox}

\end{ejercicios}


% ==============================================================================
% PESTAÑA 5: FÓRMULAS CLAVE (contentFormulas)
% ==============================================================================
% ==============================================================================
% PESTAÑA 5: FÓRMULAS CLAVE (contentFormulas)
% ==============================================================================
\begin{formulas}

  \formula{Condición Necesaria y Suficiente de Diferenciabilidad}
    {\lim\limits_{(x,y) \to (x_0, y_0)} \dfrac{f(x,y) - \left( f(x_0, y_0) + \nabla f(x_0, y_0) \cdot (x-x_0, y-y_0) \right)}{\| (x,y) - (x_0, y_0) \|} = 0}
    {Definición formal de diferenciabilidad basada en la convergencia a cero del residuo relativo a la distancia.}

  \formula{Vector Gradiente ($\nabla f$)}
    {\nabla f(x_0, y_0) = \left( \dfrac{\partial f}{\partial x}(x_0, y_0), \; \dfrac{\partial f}{\partial y}(x_0, y_0) \right)}
    {Vector de derivadas parciales de primer orden que determina la transformación lineal de la aproximación de primer orden.}

  \formula{Ecuación del Plano Tangente (Forma Extendida)}
    {z = f(x_0, y_0) + \left[ \dfrac{\partial f}{\partial x}(x_0, y_0) \right] (x - x_0) + \left[ \dfrac{\partial f}{\partial y}(x_0, y_0) \right] (y - y_0)}
    {Representación cartesiana explícita de la aproximación lineal a la superficie en el punto de evaluación.}

  \formula{Ecuación del Plano Tangente (Forma Vectorial)}
    {z = f(x_0, y_0) + \nabla f(x_0, y_0) \cdot (x - x_0, y - y_0)}
    {Representación compacta de la aproximación lineal utilizando el producto escalar entre el vector gradiente y el vector de desplazamientos coordenados.}

  \formula{Condición de Paralelismo entre Planos}
    {\vec{n}_1 = k \vec{n}_2 \quad \text{con } k \in \mathbb{R} \setminus \{0\}}
    {Dos planos son paralelos si y solo si sus vectores normales son colineales (proporcionales). Para un plano tangente, $\vec{n} = (f_x, f_y, -1)$.}

  \formula{Tangencia entre Superficies}
    {f(x_0, y_0) = g(x_0, y_0) \quad \text{y} \quad \nabla f(x_0, y_0) = \nabla g(x_0, y_0)}
    {Condición analítica para que dos superficies $z=f(x,y)$ y $z=g(x,y)$ sean tangentes: deben intersecarse en el punto y compartir exactamente el mismo vector gradiente.}

  \formula{Ecuación Segmentaria o Simétrica del Plano}
    {\dfrac{x}{a} + \dfrac{y}{b} + \dfrac{z}{c} = 1}
    {Representación algebraica de un plano que interseca los ejes coordenados en los puntos $(a,0,0)$, $(0,b,0)$ y $(0,0,c)$, asumiendo $a, b, c \neq 0$.}

\end{formulas}

\end{document}